
\documentclass[tikz,border=2pt]{standalone}
\usepackage[utf8]{inputenc}
\usepackage[T1]{fontenc}
\usepackage[french]{babel}
\usepackage{amsmath,amssymb}
\usepackage[table]{xcolor}
\usepackage{siunitx}
\usepackage[version=4]{mhchem}
\usepackage{tabularx,booktabs,array}
\usepackage{tikz}
\usetikzlibrary{arrows.meta,positioning,calc,fit,decorations.pathmorphing,patterns,shapes.geometric}
\usepackage{pgfplots}
\pgfplotsset{compat=1.18}
\definecolor{primary}{RGB}{14,116,144}
\definecolor{accent}{RGB}{16,185,129}
\definecolor{dark}{RGB}{15,23,42}
\definecolor{bglight}{RGB}{241,245,249}
\definecolor{borderlight}{RGB}{203,213,225}
\definecolor{examplepink}{RGB}{236,72,153}
\definecolor{remarkblue}{RGB}{14,165,233}
\definecolor{notepurple}{RGB}{99,102,241}
\definecolor{synthgreen}{RGB}{5,150,105}
\definecolor{synthorange}{RGB}{249,115,22}
\definecolor{synthviolet}{RGB}{79,70,229}
\definecolor{primary}{RGB}{14,116,144}
\definecolor{accent}{RGB}{16,185,129}
\definecolor{dark}{RGB}{15,23,42}
\definecolor{bglight}{RGB}{241,245,249}
\definecolor{examplepink}{RGB}{236,72,153}
\definecolor{remarkblue}{RGB}{14,165,233}
\definecolor{notepurple}{RGB}{99,102,241}
\definecolor{synthgreen}{RGB}{5,150,105}
\definecolor{synthorange}{RGB}{249,115,22}
\definecolor{synthviolet}{RGB}{79,70,229}
\begin{document}
\begin{tikzpicture}[
 box/.style={draw=primary, rounded corners=3pt, fill=bglight, minimum width=3cm, minimum height=1.2cm, align=center, font=\small\bfseries},
 arr/.style={-{Stealth[length=2.5mm]}, thick, primary}
]
 \node[box] (lin) at (0,0) {Cha\^{i}ne lin\'{e}aire\\(non cyclique,\\non ramifi\'{e}e)};
 \node[box] (ram) at (4.5,0) {Cha\^{i}ne ramifi\'{e}e\\(non cyclique,\\avec branchement)};
 \node[box] (cyc) at (9,0) {Cha\^{i}ne cyclique\\(le squelette forme\\une boucle ferm\'{e}e)};

 \node[below=0.5cm of lin, font=\small] (exlin) {
 \chemfig{-[:30]-[:-30]-[:30]-[:-30]}
 };
 \node[below=0.5cm of ram, font=\small] (exram) {
 \chemfig{-[:30](-[2])-[:-30]-[:30]}
 };
 \node[below=0.5cm of cyc, font=\small] (excyc) {
 \chemfig{*6(------)}
 };
\end{tikzpicture}
\end{document}
